\documentclass[11pt]{article}
\usepackage[margin=1in]{geometry}
\usepackage{amsmath,amssymb,booktabs,hyperref}
\newcommand{\dd}{\mathrm d}
\title{A Planck-scaled flux candidate for the vacuum--force link}
\author{Zimmerman programme: separately dated CPB-PF-2026-10-16 candidate}
\date{16 October 2026}
\begin{document}
\maketitle
\section{What this construction tests}
The preceding checkpoint showed why inserting a force coefficient directly
as a function of a sequestered vacuum counterterm changes the constraint.
Here the force instead depends on a distinct fixed conjugate flux $Q$.
A specified Planck-dependent flux energy is varied together with the
gravity auxiliaries. This produces an explicit conditional classical link:
one flux can control both a retained vacuum term and the local cubic-force
coefficient. The ordinary unscaled flux energy is a control case which
fails to survive the projected stress equation.

This is a new candidate action family for this programme, not an amendment
to the frozen theory, a derived microscopic medium or whole-theory closure.
The functions, couplings, boundary ensemble and quantum assumptions are
stated below. The absolute flux value and fitted normalization are not
predicted. Four-forms, variable gravitational couplings and self-tuning
are established research areas \cite{pf:sequester,pf:qtheory,pf:planck};
the bounded source comparison does not establish global novelty.

\section{One action and its independent variables}
Use gravity signature $(-+++)$ and natural units in this report. Write
$s=M_G^2>0$ for the gravitational multiplier, $s_*>0$ for a reference
squared Planck mass, $\Lambda$ for a vacuum counterterm auxiliary, and
$Q>0$ for a separate flux variable. The action family is
\begin{align}
 S={}&\int\dd^4x\sqrt{-g}\left[{sR\over2}-\Lambda
               -V_r(s,Q)+\mathcal L_m(g,\psi,Q)\right]\nonumber\\
 &+\int f(\Lambda/\mu^4)F_\Lambda
       +\int\widehat f(s/s_*)F_s+\int QF_Q+S_{\partial,Q},\\
 V_r(s,Q)&={Q^2\over2}(s/s_*)^r.
 \label{pf:action}
\end{align}
The three four-forms are field strengths of independent three-form gauge
potentials. Smooth nonlinear $f,\widehat f$ have nonzero derivatives on
the chosen branch. They implement the usual local constraint mechanism
of Ref.~\cite{pf:sequester}; the added $V_r$ and $Q$ sector are the
candidate extension derived here.

The boundary term $S_{\partial,Q}$ selects fixed conjugate $Q$ for the
third three-form. In an electric description it is the Legendre boundary
term, so the reduced potential is $+V_r$, rather than the sign obtained by
substituting an on-shell field strength into an unrestricted action.
The integrated $F_Q$ is not independently fixed as another boundary datum
in this ensemble. Membranes and flux jumps are absent.

The matter hypothesis is explicit: $\mathcal L_m$ is a generally covariant,
regulated matter/frame action with no explicit dependence on $s$ or
$\Lambda$, while dependence on $Q$ is allowed and varied. Its fields
would have to include the carrier, response medium and physical rest frame.
The current preferred-frame nodal Hamiltonian has not yet supplied that
complete action. The gravity calculation below is conditional on this
matter hypothesis. Renormalization of curvature operators and graviton
loops are not established by it.

Dimensions are $[s]=[s_*]=[Q]=E^2$ and $[V_r]=[\Lambda]=E^4$.
Neither $s$ nor the flux variable $Q$ is the fitted acceleration coefficient
$\kappa_a$. The flux is also unrelated to the carrier's Noether charge.

\section{Vary first; eliminate afterwards}
Gauge-potential variation gives
\begin{equation}
 \partial_\mu\Lambda=\partial_\mu s=\partial_\mu Q=0
\end{equation}
on a connected source-free branch. With
$T_{\mu\nu}=-(2/\sqrt{-g})\delta S_m/\delta g^{\mu\nu}$, the
on-shell local metric equation is
\begin{equation}
 sG_{\mu\nu}=T_{\mu\nu}-(\Lambda+V_r)g_{\mu\nu}.
 \label{pf:metric}
\end{equation}
Before imposing constancy, the usual
$(\nabla_\mu\nabla_\nu-g_{\mu\nu}\Box)s$ term is also present.
The other variations include
\begin{align}
 {f'\over\mu^4}F_\Lambda&=\mathrm{vol}_4,\\
 {\widehat f'\over s_*}F_s
      &=\left(\partial_s V_r-{R\over2}\right)\mathrm{vol}_4,\\
 F_Q&=\left(\partial_QV_r-\partial_Q\mathcal L_m\right)
                                                  \mathrm{vol}_4.
 \label{pf:auxiliary}
\end{align}
In particular a $Q$-dependent force or matter vacuum term contributes to
the last equation. It is not dropped or treated as an unchanged field
strength.

Let $\langle\cdot\rangle$ be a well-defined normalized spacetime average,
with $c_\Lambda=\int F_\Lambda$ and $c_s=\int F_s$. The first two
constraints imply
\begin{align}
 \langle R\rangle&=2\partial_sV_r+\mathcal R_F,\\
 \mathcal R_F&=-{2\mu^4\widehat f'\over s_*f'}
                                      {c_s\over c_\Lambda},\qquad
 \Delta_F={s\mathcal R_F\over4}.
\end{align}
Taking the trace of Eq.~\eqref{pf:metric}, averaging, and using this
constraint yields
\begin{align}
 \Lambda+V_r&={\langle T\rangle\over4}
                   +{s\partial_sV_r\over2}+\Delta_F,\\
 sG_{\mu\nu}&=T_{\mu\nu}-{\langle T\rangle\over4}g_{\mu\nu}
                -\left[{s\partial_sV_r\over2}+\Delta_F\right]g_{\mu\nu}.
 \label{pf:projected}
\end{align}
These equations are the result of the new action's variations; the old
projected equation with an extra flux stress appended afterwards would
not give the same answer.

For the selected power family, $s\partial_sV_r=rV_r$. Thus the retained
flux contribution to the last coefficient is $rV_r/2$.
An unscaled flux, $r=0$, disappears from that coefficient. If the design
requirement is to retain its entire positive fixed-flux energy, the power
is selected as $r=2$. This conditional selection is not a symmetry-based
derivation of that functional dependence.

Split $T_{\mu\nu}=-\epsilon_{\rm vac}(Q)g_{\mu\nu}
+\tau_{\mu\nu}$. Because $Q$ is constant, its constant matter vacuum
contribution cancels algebraically in Eq.~\eqref{pf:projected}, leaving
\begin{equation}
 sG_{\mu\nu}=\tau_{\mu\nu}-B_{\rm eff}g_{\mu\nu},\qquad
 B_{\rm eff}={rV_r\over2}+\Delta_F+{\langle\tau\rangle\over4}.
 \label{pf:Beff}
\end{equation}
The same $\epsilon_{\rm vac}'(Q)$ remains in the $F_Q$ constitutive
equation. At fixed conjugate $Q$ it changes that field strength, not the
chosen flux value by itself. This algebraic cancellation is not a proof
that a full global solution with unchanged auxiliary flux data exists
for every loop correction. The residual, boundary and matter hypotheses
remain part of the physical problem.

There is a useful check on what this construction explains. The invertible
field redefinition $\Lambda'=\Lambda+V_r(s,Q)$ removes $V_r$ from the
bulk vacuum term but changes the topological function to
$f([\Lambda'-V_r(s,Q)]/\mu^4)F_\Lambda$. Its $s,Q$ derivatives carry
exactly the same new constraints. Thus the selected power can equivalently
be regarded as a chosen cross-coupling of the four-form constraints; it
is not a parameter-free explanation hidden by a choice of variables.

\section{A fixed-flux force sector with the required scaling}
For a local static test, choose the matter response energy at fixed $Q$:
\begin{equation}
 \mathcal E_{\rm resp}={C\over2}
       |\nabla\varphi-\mathbf P/C|^2
       +W(Q,p)+q_b\rho_b\varphi,\qquad
 W(Q,p)={\alpha p^3\over Q},\quad p=|\mathbf P|,
 \quad C,\alpha>0.
 \label{pf:force}
\end{equation}
Here $[\varphi]=E$, $[p]=E^2$, $[q_b]=E^{-1}$ and $C,\alpha$ are
dimensionless. This inverse-flux functional is a new model input. It has
not been obtained from the original scalar particle action or a protective
symmetry. In particular it is not a consequence of the earlier affine
four-form mass portal, whose failure to select a halo law remains intact.

The action can retain this dependence while keeping $\mathcal L_m$
independent of $\Lambda,s$. A possible leading nodal matching would
require $A h(Q)^3=3\alpha/Q$, or
$h(Q)=[3\alpha/(AQ)]^{1/3}$. That specifies another coupling function;
it does not derive it. If the fermions are retained, their generated cubic
must not also be added as an independent $W$. Their nonlocal gradients,
temperature, occupancy and counterterms still have to be treated.

For $r=2$ the classical fixed-flux energy at fixed $s$ is
\begin{equation}
 \mathcal U(Q,p)=V_2(s,Q)+\alpha p^3/Q.
\end{equation}
Its conjugate electric field strength, before additional matter terms, is
\begin{equation}
 f_Q=Q(s/s_*)^2-\alpha p^3/Q^2,\qquad
 \partial_Qf_Q=(s/s_*)^2+2\alpha p^3/Q^3>0.
\end{equation}
The electric Legendre map is therefore monotone on the positive-$Q$
branch. The joint radial Hessian obeys
\begin{equation}
 \det\operatorname{Hess}_{Q,p}\mathcal U
 ={6\alpha p\over Q}(s/s_*)^2
             +{3\alpha^2p^4\over Q^4}>0\quad(p>0).
\end{equation}
The vector tangential curvature is $\partial_pW/p=3\alpha p/Q\geq0$.
These are classical auxiliary/static convexity results. They do not prove
full matter/metric dynamics, and an arbitrary $Q$-dependent loop vacuum
function could alter the electric Legendre curvature. The branch is
singular at $Q=0$; no smooth zero-vacuum continuation is claimed.

Variation of Eq.~\eqref{pf:force} gives
\begin{equation}
 \nabla\cdot\mathbf D=q_b\rho_b,\qquad
 \mathbf D=C\nabla\varphi-\mathbf P,\qquad
 {\mathbf D\over C}={3\alpha p\over Q}\mathbf P.
\end{equation}
Outside a spherical source, in an aligned local branch,
\begin{equation}
 D={q_bM_b\over4\pi r^2},\qquad
 p={B_P\over r},\qquad
 B_P=\sqrt{{q_bM_bQ\over12\pi C\alpha}},\qquad
 \varphi'={D+p\over C}.
\end{equation}
The scalar acceleration on baryons has magnitude $q_b\varphi'$.
The $p$ term gives a flat-speed coefficient on an annulus where it dominates:
\begin{equation}
 {v_{\rm flat}^4\over M_b}={q_b^3Q\over12\pi C^3\alpha}.
 \label{pf:TF}
\end{equation}
An infinite isolated halo is not established by this local solution.

\section{An executed one-scale nodal realization of the leading function}
There is a less arbitrary way to obtain the required leading dependence
than prescribing $h(Q)$ alone: postulate that all dimensionful parameters
of the nodal medium scale with $\mu_Q=\sqrt Q$,
\begin{equation}
 m_*=c_m\mu_Q,\quad k_0=c_k\mu_Q,\quad
 \Lambda_F=c_\Lambda\mu_Q,\quad h={c_h\over\mu_Q},\quad
 A=A_*\mu_Q,\quad A_*={Nc_m\over2\pi v}.
\end{equation}
Here $\Lambda_F$ is a band cutoff, not the gravity counterterm. This
common-scaling ansatz is additional physics whose origin remains open.
Renormalization and the other sectors can add independent scales or
logarithms; the finite-cutoff scaling test is not a proof of protection
under those corrections.
It yields
\begin{equation}
 {A h^3\over3}p^3={A_*c_h^3\over3}{p^3\over Q},\qquad
 \alpha={A_*c_h^3\over3}.
\end{equation}
The previously calculated finite-cutoff determinant, with the same critical
quadratic subtraction, scales more generally as
\begin{equation}
 V_{F,\rm ren}(Q,p)=Q^2\mathcal F(p/Q),
\end{equation}
not exactly as a cubic at every field amplitude. This follows directly
from its radial expression
\begin{equation}
 V_{F,\rm ren}=A\int_0^{\Lambda_F}
 {\Delta_f^4\,\dd t\over2[\sqrt{t^2+\Delta_f^2}+t]^2},\qquad
 \Delta_f=hp,
\end{equation}
by scaling all energies with $\mu_Q$. Its small-gap limit is the required
$\alpha p^3/Q$. The critical quadratic counterterm coefficient
$A\Lambda_Fh^2/2=A_*c_\Lambda c_h^2/2$ is independent of $Q$;
it is still a chosen tuning, not a protected symmetry result.

For a dimensionless illustration choose
$N=2$, $v=c_h=c_\Lambda=1/2$, $c_m=2\sqrt2$, $c_k=2$.
Then $A_*=4\sqrt2/\pi$ and
$\alpha=\sqrt2/(6\pi)\simeq.0750264$. These constants are chosen, not
derived; their normalization matches the calibrated example below.
The full circular nodal tube lies within its radial domain because
$c_\Lambda<c_k^2/(2c_m)=1/\sqrt2$.
The maximum radial group speed in this tube is
$\sqrt{c_k^2+2c_mc_\Lambda}/c_m=.923880$ and the axial speed is $1/2$.
This is a bounded-band group-speed check, not a proof of causal ultraviolet
front propagation or of coupled clock/metric stability.

The finite determinant was computed by independent physical radial
quadrature and a 60-digit closed expression for three $Q$ values and
$p/Q=.001,.01,.1$. After dividing by $Q^2$ the curves coincide.
Their ratios to the cubic approximation are respectively
$.999625$, $.996250$, $.962562$. Thus the claimed leading law is controlled
only in the small-gap regime; the exact finite-cutoff force still needs
its own solution.

It is also necessary to keep track of the filled reference sea. Before
the vacuum subtraction, its contribution is
\begin{equation}
 \epsilon_{\rm sea}(Q)=-{A\Lambda_F^3\over3}
       =-{A_*c_\Lambda^3\over3}Q^2.
\end{equation}
Any part identified as a genuine constant scalar-potential term cancels
in the projected metric equation, while its $Q$ derivative remains in
Eq.~\eqref{pf:auxiliary}. The homogeneous band energy alone does not prove
a vacuum-form stress: its pressure, reference inventory, frame and regulator
must be varied in a covariant completion. Other sea stresses remain in
$\tau_{\mu\nu}$ and its historic trace. With only this
additional term at $s=s_*$, empty electric curvature is
$1-2A_*c_\Lambda^3/3$. It is $.849947>0$ for the illustration.
Other cutoffs can fail that check, and other UV/matter terms are not
certified. The metric cancellation therefore must not be used to discard
the constitutive derivative or assume electric convexity automatically.

At the declared reference $V_2=(2.25\,\mathrm{meV})^4$, one has
$\mu_Q=2.67572$ meV. The original 2-MeV carrier is about
$7.47\,10^8$ times heavier. Coupling that sector to the medium introduces
a serious scale and state-matching obligation. Neither the common-scaling
ansatz, the required cold fermion state, nor its protection from the
carrier/mediator sector has been established.

\section{Which Newton constant and which vacuum energy?}
The $D$ term is an additional inverse-square force. For an unscreened
static baryon coupling it must be included once:
\begin{equation}
 G_E={1\over8\pi s},\qquad
 G_\varphi={q_b^2\over4\pi C},\qquad
 G_{\rm lab}=G_E+G_\varphi.
\end{equation}
A physical screened or relativistic matter coupling could change this
dictionary and would have to be derived. Our dimensionless example has
$G_{\rm lab}=C=q_b=1$, hence $G_E=1-1/(4\pi)>0$ and
$s=s_*=[8\pi G_E]^{-1}$. This is a bookkeeping example, not a Solar
System, fifth-force, cosmological or lensing fit. Inferring vacuum density
from a Hubble rate would also require this theory's gravity dictionary.

At the reference Planck sector $s=s_*$, the empty $r=2$ flux has
$V_2=Q^2/2$. Defining the annular acceleration scale by
$G_{\rm lab}a_0=v_{\rm flat}^4/M_b$, one obtains
\begin{align}
 a_0&={q_b^3Q\over12\pi C^3\alpha G_{\rm lab}},\\
 a_0&=\kappa_{\rm eff}\sqrt{G_{\rm lab}V_2},\qquad
 \kappa_{\rm eff}={q_b^3\over
                6\sqrt2\pi C^3\alpha G_{\rm lab}^{3/2}}.
 \label{pf:link}
\end{align}
In SI units the second line is
$a_0=\kappa_{\rm eff}c\sqrt{G_{\rm lab}\rho_4}$ with
$\rho_4=V_2/c^2$. The square-root dependence on the flux energy follows
from the one action at fixed measured Planck sector. The value
$\kappa_{\rm eff}=1/2$ would still be a calibration of the displayed
coupling combination, not a prediction. Keeping it fixed across other
$s$ sectors would impose further coupling relations and could change
the auxiliary variations.

Equation~\eqref{pf:Beff} identifies the remaining condition for replacing
$V_2$ by the actual dark-energy coefficient:
\begin{equation}
 B_{\rm eff}=V_2+\delta B,\quad
 \delta B=\Delta_F+\langle\tau\rangle/4,\qquad
 {a_0^2\over\kappa_{\rm eff}^2G_{\rm lab}B_{\rm eff}}
                   ={V_2\over V_2+\delta B}.
\end{equation}
Thus a flux normalization, auxiliary flux/history data and the full
nonvacuum historic trace remain to be fixed. The target is not obtained
by silently setting these terms to zero.
Indeed, in dimensionless reference-sector units with $\langle\tau\rangle=0$,
the allowed choices $(Q,\Delta_F)=(1,5/2)$ and $(2,1)$ both have
$B_{\rm eff}=3$, while their force scales differ by a factor two.
This is an explicit counterexample to claiming that the action with
unrestricted residual data predicts $a_0$ from the measured dark energy.
The additional boundary/history condition is a real obligation.

As a conditional example, zero $c_s$ would set $\Delta_F=0$.
For a regulated future de Sitter region with a conserved dust density,
$a(t)=e^{Ht}$ and $\tau=-\rho_{m0}e^{-3Ht}$ give
\begin{equation}
 \langle\tau\rangle_{[0,T]}
       =-{3H\rho_{m0}T\over e^{3HT}-1}\longrightarrow0.
\end{equation}
This only shows how a normalized average can vanish under stated late-time
and boundary assumptions. It is not a solution for the actual cosmology
or fixed four-form flux constraints. Infinite-volume limits, earlier
history and the occupied medium's stress require their own audit.

An empty local constant-curvature branch can be exhibited without that
averaging example. For $r=2$, $\tau_{\mu\nu}=0$, $F_s=0$, fixed
$Q,s>0$ and $\Lambda=-\epsilon_{\rm vac}(Q)$, take
\begin{equation}
 H^2={V_2\over3s},\qquad R={4V_2\over s},\qquad
 F_Q=[\partial_QV_2+\epsilon_{\rm vac}'(Q)]\mathrm{vol}_4.
\end{equation}
The metric equation and $R/2=\partial_sV_2$ are then both satisfied;
$F_\Lambda$ is fixed by its own equation. This proves consistency of
these local empty-background equations under the stated conditions. It
does not solve a global fixed-boundary problem, choose the fluxes, or
establish perturbative health of an occupied cosmological solution.
On this empty branch the SI consistency relation is
\begin{equation}
 {a_0\over cH_\Lambda}
   =\kappa_{\rm eff}\sqrt{{3G_{\rm lab}\over8\pi G_E}}.
\end{equation}
Here $H_\Lambda$ is the vacuum curvature scale, not the matter-dependent
$H(z)$. Constant $Q$ does not imply a time-varying $a_0$ proportional to
the full expansion rate through a matter era.

\section{Finite stock, energy and remaining physical tests}
The response has positive excess energy but an infinite-tail problem:
\begin{equation}
 E_W(R_1,R_2)={4\pi\alpha B_P^3\over Q}
                         \ln{R_2\over R_1}.
\end{equation}
It diverges if the local $1/r$ branch is taken to infinite radius.
A physical finite response region and matching phase are therefore
mandatory, and could change the force and projected lensing signal.
The flat coefficient above is an annular result, not a complete halo.
The field $\mathbf P$ is also not thereby identified with a conserved
displacement of the cold mass. The earlier finite-supply and mass-transport
obligations are not removed.

At fixed $Q,s$, the pure flux energy is constant. Changes in $W$ exchange
energy with the response/matter fields and sources through the common
action; they do not establish a homogeneous conversion of vacuum energy
into cold particles. The original pressure-aware exchange null relation
remains a separate consistency condition. No membrane-driven $Q$ history,
particle abundance or charge asymmetry has been generated here.

This candidate supplies a conditional classical gravity/force link and
retains the $Q$-dependent matter variation instead of confusing $Q$ with
the vacuum counterterm. It still needs a microscopic origin and radiative
control of $\alpha p^3/Q$, a covariant occupied-band/frame action, actual
cold production/transport, a physical interface and energy receiver,
metric/lensing and local-gravity tests, native cosmology and joint
empirical validation. The absolute flux, residual boundary data and
$\kappa$ are not derived. No whole-theory or global-novelty claim is made.

\section{Evidence and literature scope}
\texttt{planck\_flux\_link\_run\_002} records 63 bounded controls:
exact trace/auxiliary algebra, monotone electric inversions, sample static
Hessians, force roots, calibrated normalization, logarithmic energy and a
late-average example. The general identities are proved in the equations
above, not by finite sampling. Run 001 is the earlier completed calculation;
run 002 explicitly adds the Newtonian constant dictionary. This is self-review.
The finite-cutoff one-scale test is
\texttt{one\_scale\_nodal\_flux\_run\_002}, with 17 controls.
Its preceding completed run used the same mathematics; run 002 adds the
explicit mpmath version to the execution metadata.

The local sequestering ancestor was checked at its precise field equations.
Primary abstracts for Refs.~\cite{pf:qtheory,pf:planck} establish related
four-form/gravitational-coupling and Planck-variation research. Their detailed
solutions are not imported here. The searches are bounded context/overlap
checks; neither absent matches nor a different notation establishes novelty.

\begin{thebibliography}{9}
\bibitem{pf:sequester} N. Kaloper, A. Padilla, D. Stefanyszyn and
G. Zahariade, A Manifestly Local Theory of Vacuum Energy Sequestering,
\href{https://arxiv.org/pdf/1505.01492v2}{arXiv:1505.01492v2},
Eqs.~(7)--(10). The new $V_r$ and third flux must be varied separately.
\bibitem{pf:qtheory} F. R. Klinkhamer and G. E. Volovik,
Dynamic vacuum variable and equilibrium approach in cosmology,
\href{https://arxiv.org/abs/0806.2805v7}{arXiv:0806.2805v7}.
Its four-form amplitude is not identified with our conjugate $Q$.
\bibitem{pf:planck} D. Sobral-Blanco and L. Lombriser,
Exploring the self-tuning of the cosmological constant from Planck mass
variation, \href{https://arxiv.org/abs/2012.01838v3}{arXiv:2012.01838v3}.
\end{thebibliography}
\end{document}
